\documentclass[10pt,a4paper]{amsart}
\usepackage[margin=19mm]{geometry}
\usepackage{amsmath,amssymb,mathtools}
\usepackage[scr=rsfs,cal=euler]{mathalfa}
\usepackage{xfrac}
\usepackage[unicode,hidelinks,bookmarks=false]{hyperref}
\newcommand{\ie}{i.e.\ }
\newcommand{\cf}{cf.\ }
\newcommand{\resp}{respectively\ }
\newcommand{\Subsection}{\subsection}
\providecommand{\nobreakdash}{\nobreak\hbox{-}}
\setlength{\emergencystretch}{2em}
\multlinegap0pt
\usepackage{amssymb}
\usepackage{enumerate}
\usepackage{tikz}
\usepackage{cancel}
\usepackage{mathtools}
\usepackage{xcolor}
\newtheorem{remark}{Remark}
\newtheorem{corollary}{Corollary}
\def\depth{\operatorname{depth}}
\title{On the existence of the maximal ideal in the set of associated primes of monomial ideals}
\author{M. Cimpoea\c{s}, M. Nasernejad, A. A. Qureshi}
\date{Source: arXiv:2608.28243v1 (CC BY 4.0)}
\begin{document}
\maketitle

\noindent\emph{Source and licence.} On the existence of the maximal ideal in the set of associated primes of monomial ideals. Version arXiv:2608.28243v1 (CC BY 4.0), distributed by arXiv under the Creative Commons Attribution 4.0 International licence, http://creativecommons.org/licenses/by/4.0/, as recorded in the arXiv licence metadata for this exact version and shown on the abstract page https://arxiv.org/abs/2608.28243v1. Full licence notice: \texttt{licenses/arxiv-2608.28243-CC-BY-4.0.txt}.\par\smallskip

\noindent\textit{Editorial scope.} Counted unit: the article's corollary C4 (source0.tex lines 902--920). The article's complete original proof of it is delivered with it (source0.tex lines 922--957, one \texttt{proof} environment of 36 lines). No mathematics is written, repaired or re-derived: the statement and the proof are the article's own lines, in the article's own notation, and no retained line is substituted or re-flowed.  Retained uncounted as the source-local context the proof needs: lines 469--471; lines 804--815.  Nothing was omitted from any retained span, so the package carries no omission record.  No retained line contains a citation, so the wrapper carries no bibliography and no bibliography entry.  No retained line contains a figure environment, an image-inclusion command or drawing code, so no image is delivered and the package declares no asset dependency.  Editorial formatting only, in addition: an AMS \texttt{amsart} preamble that restates the article's own declarations and macros used by the retained lines, namely the environments \texttt{remark}, \texttt{corollary} on separate counters, together with the standard packages those lines need.  The retained environments print in this document as Remark 1, Corollary 1 in order of appearance, whereas the article prints the same items with the numbering of its own full document.  The complete native source of the article as distributed in the arXiv e-print for this exact version is bundled unchanged as \texttt{sources/208772/source0.tex}.
\begin{remark}\label{rem1}
Let $I\subset R=K[x_1,\ldots,x_n]$ be a proper  nonzero  monomial ideal. Then $\depth(R/I)=n-1$ if and only if $I$ is principal and, otherwise, $\depth(R/I)\leq n-2$.
\end{remark}

\begin{corollary}\label{c2}
Let $p\geq 1$ and $i_p>i_{p-1}>\cdots>i_1>i_0=0$ be some integers. Let $I_j\subset K[y,z]$, $0\leq j\leq p$, be some nonzero 
proper monomial ideals. We consider the monomial ideal $$I=x^{i_p}I_p + \cdots + x^{i_1}I_1 + I_0 \subset R=K[x,y,z],$$
 and we assume that
$\mathcal{G}(I)=\mathcal{G}(x^{i_p}I_p) \cup \cdots \cup \mathcal{G}(x^{i_1}I_1) \cup \mathcal{G}(I_0).$ 
Then the following statements are  equivalent:
\begin{enumerate}
\item[(1)] $\depth(R/I)=0$.
\item[(2)] There exists $0\leq q<p$ and a monomial $u\in I_{q+1}$, $u\notin I_j$  for $0\leq j\leq q$, 
           such that $|\mathcal{G}((I_0:u)+\cdots+(I_q:u))|\geq 2$.
\end{enumerate}
\end{corollary}

\begin{corollary} \label{C4}
Let $p\geq 1$ and $i_p>i_{p-1}>\cdots>i_1>i_0=0$ be some integers. Let $u_j\in K[y,z]$, $0\leq j\leq p-1$, be some monomials such that
there exists $0\leq q\leq p-1$ with $u_q\mid u_{q-1} \mid  \cdots \mid u_1 \mid u_0$, the divisions are strict, and,
if $q\leq p-2$, then
\begin{enumerate}
\item[(i)]   $u_j\nmid u_k$ for all $p-1\geq k>j\geq q$,
\item[(ii)]  $((u_{q},u_{q+1},\ldots,u_{j-1}):u_j)$ is principal  for all $q+1\leq j\leq p-1$,
\item[(iii)] $((u_{q},u_{q+1},\ldots, u_{p-1}):u)$ is principal  for any $u\in \mathcal G(I_p)$.
\end{enumerate}
Let $I_p\subset K[y,z]$ be a proper nonzero monomial ideal such that $u_q,\ldots,u_{p-1}\notin I_p$. We consider the ideal
$$I=x^{i_p}I_p+(x^{i_{p-1}}u_{p-1},\ldots,x^{i_1}u_1,u_0)\subset R=K[x,y,z].$$
We assume that \[
\mathcal{G}(I)
=
\mathcal{G}(x^{i_p}I_p)\cup\cdots\cup
\mathcal{G}(x^{i_1}I_1)\cup\mathcal{G}(I_0).
\]
Then $\depth(R/I)=1$.
\end{corollary}

\begin{proof}
%Let $I_j=(u_j)$  for $0\leq j\leq p-1$. We claim that 
%$$\mathcal G(I)=\mathcal G(x^{i_p}I_p) \cup \cdots \cup \mathcal G(x^{i_1}I_1) \cup \mathcal G(x^{i_0}I_0). $$
%Indeed, the strict inequalities $i_p > i_{p-1} > \cdots > i_1 > i_0 = 0$ prevent any generator
%involving a higher power of $x$ from dividing a generator involving a
%lower power of $x$. The hypotheses  $u_q\mid u_{q-1} \mid  \cdots \mid u_1 \mid u_0$, and the divisions are strict, 
% $u_j\nmid u_k$ for all $p-1\geq k>j\geq q$, and $u_q,\ldots,u_{p-1}\notin I_p$, guarantee the converse. 
For $0\leq j\leq p-1$, let $L_j=I_0+\cdots+I_j$. Since $u_q\mid u_{q-1} \mid  \cdots \mid u_1 \mid u_0$, we get
$$I_0\subset I_1 \subset \cdots \subset I_q,$$
and, consequently, $L_j=I_j$ for $0\leq j\leq q$. By Remark \ref{rem1}, we obtain $\depth(R'/L_j)=1>0$ for all $0\leq j\leq q$.

Now, assume that $q\leq p-2$. Note that, for any $q+1\leq j\leq p-1$, we have the following equalities 
$$L_j=I_0+\cdots+I_j=I_q+\cdots+I_j=(u_q,u_{q+1},\ldots,u_j).$$ 
If $L_j$ is not principal, then, according to Remark \ref{rem1}, $\depth(R'/L_j)=0$. 
Since $L_q = I_q = (u_q)$  is  principal, we can deduce that $(L_q:u)$ is principal as well. 
If $q+2\leq j\leq p-1$,   then, for any monomial $u\in I_{j}\setminus L_{j-1}$, we claim that
$(L_{j-1}:u)$ is principal. Indeed, since $u\in I_{j}$, it follows that $u=u_{j}u'$ for some monomial $u'$.
Due to  (ii), the ideal $(L_{j-1}:u_{j}) = ((u_{q},u_{q+1},\ldots,u_{j-1}):u_j)$ is principal, 
and, therefore, $(L_{j-1}:u)=((L_{j-1}:u_j):u')$ is principal.
Similarly, according to (iii), for any monomial $u\in I_p\setminus L_{p-1}$, we have  $(L_{p-1}:u)$ is principal.
To apply Corollary \ref{c2}, we need to verify that the condition
corresponding to statement {\rm (2)} of Corollary \ref{c2} does not occur.
 The proof above shows that, for $q+2\leq j\leq p-1,$ we have
\[
(L_{j-1}:u)\ \text{is principal for every }u\in I_{j}\setminus L_{j-1},
\]
and  for \(j=p\), $(L_{p-1}:u)$ is principal for every $u\in I_p\setminus L_{p-1}.$ 
Therefore, there is no integer \(q\) and no monomial \(u\) satisfying the
condition in statement {\rm (2)} of Corollary \ref{c2}. Hence, $\depth(R/I)\neq 0.$ 
Since $$\mathcal{G}(I)=\mathcal{G}(x^{i_p}I_p)\cup
\mathcal{G}(x^{i_{p-1}}I_{p-1})\cup\cdots\cup\mathcal{G}(I_0),$$
and \(I_p\neq 0\), the ideal \(I\) has at least two minimal generators.
Hence \(I\) is not principal. By Remark \ref{rem1},
 $\depth(R/I)\leq 1.$ Combining this with \(\depth(R/I)\neq0\), we conclude that
$\depth(R/I)=1.$
\end{proof}

\end{document}
